% ============================================================
%  PULSE Framework — Full System Architecture Document
%  People-centric Framework for Correspondence (PFC)
%  Singapore Government Digital Inclusion Platform
% ============================================================
\documentclass[12pt,a4paper]{article}

% ── Packages ─────────────────────────────────────────────────────────────────
\usepackage[T1]{fontenc}
\usepackage[utf8]{inputenc}
\usepackage[margin=2.5cm]{geometry}
\usepackage{lmodern}
\usepackage{microtype}
\usepackage{hyperref}
\usepackage{xcolor}
\usepackage{graphicx}
\usepackage{float}
\usepackage{longtable}
\usepackage{booktabs}
\usepackage{array}
\usepackage{tabularx}
\usepackage{multirow}
\usepackage{listings}
\usepackage{minted}
\usepackage{amsmath}
\usepackage{amssymb}
\usepackage{tikz}
\usepackage{pgf-umlsd}
\usepackage{enumitem}
\usepackage{fancyhdr}
\usepackage{titlesec}
\usepackage{tocloft}
\usepackage{mdframed}
\usepackage{caption}
\usepackage{subcaption}
\usepackage{fontawesome5}

% ── Colour palette ───────────────────────────────────────────────────────────
\definecolor{pulseteal}{HTML}{0b6160}
\definecolor{pulsecream}{HTML}{fffcf5}
\definecolor{pulsegold}{HTML}{d4a017}
\definecolor{pulsered}{HTML}{c0392b}
\definecolor{pulseblue}{HTML}{2c6fad}
\definecolor{codebg}{HTML}{f4f4f4}
\definecolor{bordercolor}{HTML}{cccccc}

% ── Hyperref ─────────────────────────────────────────────────────────────────
\hypersetup{
  colorlinks=true,
  linkcolor=pulseteal,
  urlcolor=pulseblue,
  citecolor=pulseteal,
  pdftitle={PULSE Framework — System Architecture},
  pdfauthor={PULSE Engineering Team},
  pdfsubject={Architecture, AI, Security, Machine Learning},
}

% ── Header / footer ──────────────────────────────────────────────────────────
\pagestyle{fancy}
\fancyhf{}
\fancyhead[L]{\small\textcolor{pulseteal}{\textbf{PULSE Framework}}}
\fancyhead[R]{\small\textcolor{gray}{System Architecture — Confidential}}
\fancyfoot[C]{\thepage}
\renewcommand{\headrulewidth}{0.4pt}

% ── Section formatting ───────────────────────────────────────────────────────
\titleformat{\section}
  {\large\bfseries\color{pulseteal}}{\thesection}{1em}{}[\titlerule]
\titleformat{\subsection}
  {\normalsize\bfseries\color{pulseteal}}{\thesubsection}{1em}{}
\titleformat{\subsubsection}
  {\normalsize\bfseries}{\thesubsubsection}{1em}{}

% ── Code listings ────────────────────────────────────────────────────────────
\lstset{
  backgroundcolor=\color{codebg},
  basicstyle=\ttfamily\small,
  breaklines=true,
  frame=single,
  rulecolor=\color{bordercolor},
  xleftmargin=6pt,
  xrightmargin=6pt,
  aboveskip=8pt,
  belowskip=8pt,
}

% ── Callout box ──────────────────────────────────────────────────────────────
\newmdenv[
  backgroundcolor=pulsecream,
  linecolor=pulseteal,
  linewidth=1pt,
  leftmargin=0pt,
  rightmargin=0pt,
  innerleftmargin=10pt,
  innerrightmargin=10pt,
  innertopmargin=6pt,
  innerbottommargin=6pt,
]{callout}

% ── Title ─────────────────────────────────────────────────────────────────────
\title{
  \vspace{-1cm}
  {\color{pulseteal}\rule{\linewidth}{2pt}}\\[0.4cm]
  {\Huge\bfseries\color{pulseteal} PULSE Framework}\\[0.2cm]
  {\Large People-centric Framework for Correspondence}\\[0.2cm]
  {\large\textit{Full System Architecture — Security, Protocols, AI/ML Pipeline}}\\[0.3cm]
  {\color{pulseteal}\rule{\linewidth}{1pt}}\\[0.3cm]
  {\normalsize Singapore Government Digital Inclusion Platform}
}
\author{PULSE Engineering Team}
\date{July 2026 — Version 3.0 (Live Production)}

% ═════════════════════════════════════════════════════════════════════════════
\begin{document}
% ═════════════════════════════════════════════════════════════════════════════

\maketitle
\thispagestyle{fancy}

\begin{callout}
\textbf{Classification:} Internal / Engineering Reference \\
\textbf{Status:} Live production system — Telegram bot active, WhatsApp Meta integration wired \\
\textbf{Repository:} \texttt{OriginalNneo/Project-PULSE} (branch: \texttt{main}) \\
\textbf{Live endpoint:} \texttt{https://pulse.nathanielbuilds.cc}
\end{callout}

\vspace{1em}
\tableofcontents
\newpage

% ─────────────────────────────────────────────────────────────────────────────
\section{Executive Overview}
% ─────────────────────────────────────────────────────────────────────────────

PULSE (People-centric Framework for Correspondence) is a multi-channel AI-mediated
correspondence platform designed to eliminate the digital accessibility gap for
Singapore's most vulnerable citizens — seniors 65+, persons with disabilities, and
individuals with low digital literacy. The system bridges CPF (Central Provident Fund)
services and citizens through empathetic, adaptive, multilingual AI agents.

\begin{callout}
\textbf{Scope of this document:} The live production system comprising a Node.js/Express
TypeScript backend, a Next.js 14 frontend, an AI multi-agent pipeline (OpenClaw), HuggingFace
Inference API models, a z.ai GLM-5-Turbo LLM (Hermes), MongoDB Atlas and SQLite persistence,
a Telegram Bot and WhatsApp Meta Cloud API channel layer, and a Caddy reverse proxy serving
\texttt{pulse.nathanielbuilds.cc}. Aspirational infrastructure (Kafka, Redis, PostgreSQL,
Kubernetes) is documented in Section~\ref{sec:aspirational} and has not yet been provisioned.
\end{callout}

\subsection{PULSE Acronym}

\begin{tabular}{@{}lll@{}}
\toprule
\textbf{Letter} & \textbf{Pillar} & \textbf{Design Intent} \\
\midrule
\textbf{P} & People-Centric   & Adaptive to each individual, not demographic buckets \\
\textbf{U} & Universal        & WCAG 2.2 AA by default, all channels, all languages \\
\textbf{L} & Linked           & Seamless analog-to-digital bridge; physical mail + voice + digital \\
\textbf{S} & Secure \& Supportive & Scam-resistant by design; empowering without exposing \\
\textbf{E} & Empathetic        & Keeps a finger on the community's emotional pulse \\
\bottomrule
\end{tabular}

\subsection{Target User Populations}

\begin{tabular}{@{}llr@{}}
\toprule
\textbf{Group} & \textbf{Primary Barrier} & \textbf{Singapore Population} \\
\midrule
Seniors (65+)             & Low digital literacy, scam anxiety, declining vision & $\sim$1 in 4 by 2030 \\
Persons with Disabilities & Screen-reader incompatibility, cognitive overload & $\sim$3.4\% of residents \\
Low Digital Literacy      & Language barriers, unfamiliarity with gov portals & Cross-demographic \\
Caregivers / Proxies      & Need delegated access with audit trail & Significant family-care market \\
\bottomrule
\end{tabular}

% ─────────────────────────────────────────────────────────────────────────────
\section{High-Level System Architecture}
\label{sec:highlevel}
% ─────────────────────────────────────────────────────────────────────────────

The production system is a monorepo with two deployed processes on a single VPS,
fronted by a Caddy reverse proxy. The overall concentric-ring security model governs
both the live production tier and the planned enterprise expansion.

\subsection{Concentric-Ring Security Model}

\begin{lstlisting}
┌───────────────────────────────────────────────────────────────────────┐
│  INTERNET / CLIENTS                                                   │
│  (Telegram users, WhatsApp users, browser users, Caddy)              │
└────────────────────────────────┬──────────────────────────────────────┘
                                 │
                                 ▼
┌───────────────────────────────────────────────────────────────────────┐
│  RING 1 — Edge & Security Perimeter                                  │
│  Caddy (TLS/HTTPS), HSTS, security headers (Helmet.js),              │
│  CORS enforcement, HMAC-SHA256 webhook verification                   │
└────────────────────────────────┬──────────────────────────────────────┘
                                 │
                                 ▼
┌───────────────────────────────────────────────────────────────────────┐
│  RING 2 — API Gateway (Express.js :3000)                             │
│  JWT/cookie auth, role-based access, rate limiting,                   │
│  request ID tracing, Zod schema validation                            │
└────────────────────────────────┬──────────────────────────────────────┘
                                 │
                                 ▼
┌───────────────────────────────────────────────────────────────────────┐
│  RING 3 — Application Services (Modular Core)                        │
│  9 service modules + AI agent system (OpenClaw)                       │
│  Inbound processor, session manager, escalation pipeline              │
└────────────────────────────────┬──────────────────────────────────────┘
                                 │
                                 ▼
┌───────────────────────────────────────────────────────────────────────┐
│  RING 4 — Data & Integration Layer                                   │
│  SQLite (customers, sessions), MongoDB Atlas (knowledge, queue),      │
│  HuggingFace Inference API, z.ai GLM Hermes, Telegram Bot API,        │
│  WhatsApp Meta Cloud API                                               │
└───────────────────────────────────────────────────────────────────────┘
\end{lstlisting}

\subsection{Live Production Deployment}

\begin{lstlisting}
                   DNS: pulse.nathanielbuilds.cc
                              │
                              ▼
                    ┌─────────────────┐
                    │  Caddy (HTTPS)  │  TLS auto-managed (Let's Encrypt)
                    │  :443 → :80     │  HTTP/2, HSTS enforced
                    └────────┬────────┘
                             │  Route rules:
           ┌─────────────────┼────────────────────┐
           │                 │                    │
     /dashboard/*       /webhook/*           /api/*, /query,
     /api/*, etc.       /dashboard/ws         everything else
           │                 │                    │
           ▼                 ▼                    ▼
  ┌──────────────┐   ┌──────────────┐   ┌──────────────┐
  │  Backend     │   │  Backend     │   │  Frontend    │
  │  Express     │   │  Express     │   │  Next.js 14  │
  │  :3000       │   │  :3000       │   │  :3001       │
  └──────────────┘   └──────────────┘   └──────────────┘
  pm2: pulse-backend                    pm2: pulse-frontend
\end{lstlisting}

Both processes are managed by \textbf{pm2} running from \texttt{/nat/Project-PULSE}.
The backend runs via \texttt{npx tsx src/gateway/index.ts} (TypeScript executed
directly by tsx; no compile step). The frontend runs via \texttt{npx next start -p 3001}.
Auto-deploy is triggered by GitHub webhook \texttt{POST /webhook/github} on pushes
to \texttt{main}.

% ─────────────────────────────────────────────────────────────────────────────
\section{Ring 1 — Edge, TLS, and Security Headers}
\label{sec:ring1}
% ─────────────────────────────────────────────────────────────────────────────

\subsection{Caddy Reverse Proxy}

Caddy handles TLS certificate provisioning and renewal automatically via ACME (Let's Encrypt).
It enforces:
\begin{itemize}[noitemsep]
  \item \textbf{HTTPS-only}: all HTTP redirected to HTTPS (301)
  \item \textbf{HTTP/2}: multiplexed, head-of-line-blocking eliminated
  \item \textbf{WebSocket upgrades}: \texttt{/dashboard/ws} proxied with \texttt{Upgrade}/\texttt{Connection}
    headers preserved
  \item \textbf{Route isolation}: \texttt{/webhook/*} endpoints not exposed to the public frontend path
\end{itemize}

\subsection{HTTP Security Headers (Helmet.js)}

The Express gateway applies \texttt{helmet()} with the following headers on every response:

\begin{tabular}{@{}lp{9cm}@{}}
\toprule
\textbf{Header} & \textbf{Value / Policy} \\
\midrule
\texttt{Strict-Transport-Security} & \texttt{max-age=31536000; includeSubDomains} — forces HTTPS for 1 year \\
\texttt{Content-Security-Policy}   & \texttt{default-src 'self'; script-src 'self'} — blocks inline scripts/XSS \\
\texttt{X-Content-Type-Options}    & \texttt{nosniff} — prevents MIME sniffing attacks \\
\texttt{X-Frame-Options}           & \texttt{DENY} — blocks clickjacking via iframes \\
\texttt{Referrer-Policy}           & \texttt{strict-origin-when-cross-origin} \\
\texttt{X-DNS-Prefetch-Control}    & \texttt{off} \\
\texttt{Permissions-Policy}        & Restricts camera, microphone, geolocation to \texttt{'none'} \\
\bottomrule
\end{tabular}

\subsection{CORS Policy}

CORS is enforced at the Express middleware layer. The allowed origins are:

\begin{lstlisting}[language=]
Allowed Origins:
  - http://localhost:3001    (local dev frontend)
  - https://pulse.nathanielbuilds.cc  (production)

Allowed Methods: GET, POST, PUT, PATCH, DELETE, OPTIONS
Allowed Headers: Content-Type, Authorization, X-Request-Id
Credentials: true (required for httpOnly cookie transmission)
\end{lstlisting}

Same-origin requests from the production domain do not trigger CORS preflight —
the browser omits the \texttt{Origin} header when the scheme, host, and port match.

\subsection{Webhook Authentication — HMAC-SHA256}

Both inbound webhook channels use HMAC-SHA256 for payload authentication.

\subsubsection{Telegram Bot API}

Telegram does not send an HMAC signature by default. The PULSE backend validates
Telegram messages by checking the \texttt{X-Telegram-Bot-Api-Secret-Token} header
(optional token set on the webhook registration). The primary defence is that the
webhook URL is not publicly advertised and the endpoint returns \texttt{200 OK}
to all requests (Telegram requires this to avoid retries); the token check is a
secondary gate.

\subsubsection{Meta WhatsApp Cloud API}

The Meta webhook handler at \texttt{POST /webhook/whatsapp} enforces:

\begin{lstlisting}[language=JavaScript]
// Signature format: "sha256=<hex-digest>"
const expected = "sha256=" +
  crypto.createHmac("sha256", process.env.WHATSAPP_APP_SECRET)
        .update(rawBody)
        .digest("hex");
crypto.timingSafeEqual(Buffer.from(expected), Buffer.from(header));
// timingSafeEqual prevents timing-oracle attacks
\end{lstlisting}

The \texttt{rawBody} buffer is preserved before JSON parsing via Express
\texttt{verify} callback, because HMAC must be computed on the raw bytes,
not the parsed object. An additional \texttt{appsecret\_proof} parameter is
appended to all outbound Graph API calls:

\[
\texttt{appsecret\_proof} = \text{HMAC-SHA256}(\texttt{WHATSAPP\_APP\_SECRET},\ \texttt{WHATSAPP\_ACCESS\_TOKEN})
\]

\subsubsection{GitHub Webhook (Auto-Deploy)}

The GitHub webhook at \texttt{POST /webhook/github} uses the same HMAC-SHA256
pattern against \texttt{GITHUB\_WEBHOOK\_SECRET}. The handler:
\begin{enumerate}[noitemsep]
  \item Verifies signature before any payload processing
  \item Filters to \texttt{push} events on \texttt{refs/heads/main} only
  \item Responds \texttt{202 Accepted} immediately (before deploy runs)
  \item Executes \texttt{deploy/deploy-frontend.sh} asynchronously via \texttt{execFile}
\end{enumerate}

% ─────────────────────────────────────────────────────────────────────────────
\section{Ring 2 — API Gateway}
\label{sec:ring2}
% ─────────────────────────────────────────────────────────────────────────────

\subsection{Express Gateway Entry Point (\texttt{src/gateway/index.ts})}

The gateway wires the following middleware stack in order:

\begin{enumerate}[noitemsep]
  \item \textbf{Helmet} — security headers
  \item \textbf{CORS} — origin/method enforcement
  \item \textbf{Cookie-parser} — parses httpOnly session cookies
  \item \textbf{Raw body preservation} — \texttt{verify} callback captures \texttt{rawBody: Buffer} before JSON parse (required for HMAC)
  \item \textbf{JSON body parser} — \texttt{express.json(\{limit: "10mb"\})}
  \item \textbf{Request ID tracing} — generates \texttt{X-Request-Id} and attaches to \texttt{req}
  \item \textbf{Rate limiter} — in-memory per-IP, role-aware
  \item \textbf{Route mounting} — see below
  \item \textbf{Global error handler} — catches \texttt{AppError} subclasses, returns structured JSON
\end{enumerate}

\subsection{Route Mounting Table}

\begin{tabular}{@{}llp{6cm}@{}}
\toprule
\textbf{Mount Path} & \textbf{Module} & \textbf{Description} \\
\midrule
\texttt{/query}              & inbound.ts           & Main citizen query endpoint (text) \\
\texttt{/query/voice}        & inbound.ts           & Voice/audio query endpoint \\
\texttt{/dashboard}          & dashboard.ts         & Officer dashboard REST API \\
\texttt{/dashboard/ws}       & ws.ts                & WebSocket real-time push \\
\texttt{/webhook/telegram}   & webhook.ts           & Telegram Bot update handler \\
\texttt{/webhook/whatsapp}   & webhook.ts           & Meta WhatsApp Cloud handler \\
\texttt{/webhook/github}     & deploy.ts            & GitHub push auto-deploy \\
\texttt{/api/v1/agents}      & orchestrator/routes  & OpenClaw agent HTTP API \\
\texttt{/api/v1/correspond.} & correspondence/routes & Service: letters, templates \\
\texttt{/api/v1/vuln}        & vulnerability/routes  & Service: vulnerability tiers \\
\texttt{/api/v1/delivery}    & delivery/routes      & Service: multi-channel dispatch \\
\texttt{/api/v1/proxy}       & proxy/routes         & Service: delegated access \\
\texttt{/api/v1/analytics}   & analytics/routes     & Service: engagement metrics \\
\texttt{/api/v1/billing}     & billing/routes       & Service: usage metering \\
\bottomrule
\end{tabular}

\subsection{Authentication and JWT}

\begin{lstlisting}
┌────────────┐    ┌──────────────┐    ┌──────────────────┐
│  Client    │───▶│  Auth Service │───▶│  Singpass/       │
│  Browser   │    │              │    │  Corppass IdP     │
└────────────┘    └──────┬───────┘    └──────────────────┘
                         │ Issue JWT pair
                         ▼
                   httpOnly cookies:
                   ┌────────────────────────────────────┐
                   │  access_token  (JWT, 15 min TTL)  │
                   │  refresh_token (opaque, 7 day TTL) │
                   └────────────────────────────────────┘

JWT payload:
{
  "sub":       "<user_id>",
  "tenant_id": "<tenant>",
  "roles":     ["user" | "proxy" | "officer" | "admin"],
  "lang":      "en" | "zh" | "ms" | "ta",
  "vuln_tier": "self-service" | "guided" | "high-touch",
  "iat":       <unix_ts>,
  "exp":       <unix_ts>,
  "jti":       "<unique token id>"  // for revocation
}
\end{lstlisting}

\textbf{Token revocation:} Revoked JTI values are stored in Redis with TTL equal to
token expiry. The gateway checks the revocation list on every authenticated request.
Refresh tokens are rotated on every use — token theft is detectable within one use.

\subsection{Rate Limiting}

\begin{tabular}{@{}lll@{}}
\toprule
\textbf{Role} & \textbf{Limit} & \textbf{Storage} \\
\midrule
Anonymous       & 10 req/min    & In-memory (per IP)      \\
Authenticated   & 100 req/min   & In-memory (per user ID) \\
Admin           & 500 req/min   & In-memory (per user ID) \\
Service account & 1000 req/min  & In-memory (per service) \\
\bottomrule
\end{tabular}

Exceeded limits return \texttt{429 Too Many Requests} with a \texttt{Retry-After}
header. The Telegram and WhatsApp webhook paths bypass rate limiting (handled
at channel level by their respective providers).

\subsection{Request Tracing}

Every inbound request receives a unique \texttt{X-Request-Id} (UUIDv4). This ID is:
\begin{itemize}[noitemsep]
  \item Attached to the Pino structured logger context for every log line
  \item Propagated in \texttt{X-Trace-Id} / \texttt{X-Span-Id} headers to downstream calls
  \item Returned in the response for client-side error correlation
\end{itemize}

% ─────────────────────────────────────────────────────────────────────────────
\section{Ring 3 — Inbound Message Processing Pipeline}
\label{sec:inbound}
% ─────────────────────────────────────────────────────────────────────────────

The inbound pipeline is the heart of the live system. It processes every message from
a Telegram user or WhatsApp user through a multi-stage AI pipeline before delivering
a reply.

\subsection{Full Processing Pipeline}

\begin{lstlisting}
Telegram update / WhatsApp webhook
  │
  ▼
[1] CHANNEL ADAPTER
    telegram.ts / webhook.ts
    - Validates update type (message | callback_query | voice)
    - Extracts userId (tg:<chatId> | wa:<phone>), text, voice OGG
    - Handles inline keyboard callbacks (escalation button, CSAT stars)
    │
    ▼
[2] INBOUND ROUTER (processInbound)
    src/gateway/inbound.ts
    │
    ├── handleCommand()        /end, /start, /help — fast exit
    ├── runStateIntercepts()   CSAT rating collection, guided Q&A, pending offer
    │
    ▼
[3] PARALLEL PRE-PROCESSING (Promise.all)
    ├── Emotion Detection   ──▶  HF: j-hartmann/emotion-english-distilroberta-base
    ├── Language Detection  ──▶  [script] → [HF: xlm-roberta] → [GLM fallback]
    └── Voice Processing    ──▶  (if voice) HF: openai/whisper-large-v3 + duration
    │
    ▼
[4] TRIAGE CLASSIFICATION
    src/agents/triage/classifier.ts
    Cat 3: personal account data (regex)
    Cat 2: guiding-question topics (regex)
    Cat 1: public knowledge (default)
    │
    ▼
[5] VOICE CLARITY CHECK (voice only)
    - durationSec < 1.2s → reject
    - text.length < 5     → reject
    - text.length / durationSec < 4 → NOISY ENVIRONMENT path
    - transcriptionLooksUnclear() → GARBLED path
    │
    ▼
[6] GUIDING QUESTIONS (Cat 2)
    MongoDB cpf_guiding_questions collection
    One question asked at a time; synthesises answer after final Q
    │
    ▼
[7] QUERY AGENT (Cat 1 / Cat 3 / completed guided)
    Emotion trajectory folding → effectiveEmotion
    Tone directive built → toneDirective(score, label, sustained)
    RAG retrieval: MongoDB cpf_knowledge
    GLM-5-Turbo generation (Hermes, z.ai)
    Response formatted → formatter.ts
    │
    ▼
[8] ESCALATION ANALYSIS
    analyzeEscalation(query, reply, confidence, triage)
    5 layers: personal_data | life_event | explicit_request | complaint | bot_failure
    emotion_score > 70 → auto-escalate
    │
    ▼
[9] DELIVER ANSWER
    sendReply() → Telegram sendMessage / WhatsApp messages API
    If escalated → postToQueue() → MongoDB ccu_queue
    If !langConfident && firstTurn && !English → append soft language hint
    │
    ▼
[10] ASYNC POST-PROCESSING
    updateQueuePriority() — enriched score (financial + sustained)
    pushEmotionEvent()   — dashboard ring buffer
    broadcast() WS       — live dashboard push
\end{lstlisting}

\subsection{Voice Processing Sub-pipeline}

When a Telegram user sends a voice note (OGG/Opus):

\begin{enumerate}[noitemsep]
  \item Telegram API: download OGG file via \texttt{getFile} + \texttt{file\_path}
  \item HuggingFace \texttt{openai/whisper-large-v3}: speech-to-text transcription
  \item HuggingFace \texttt{audeering/wav2vec2-large-robust-12-ft-emotion-msp-dim}: dimensional audio emotion
  \item Duration extracted from OGG header (used for noise ratio check)
  \item Background noise check: $\texttt{len(text)} / \texttt{durationSec} < 4$
\end{enumerate}

\subsection{Thinking Animation}

During processing, the bot sends a \texttt{💭} typing indicator via
\texttt{sendChatAction("typing")} with 2-second intervals (Telegram's
rate limit). The indicator is a best-effort, non-blocking async process;
failures are caught and logged at \texttt{warn} level and do not affect the reply.

% ─────────────────────────────────────────────────────────────────────────────
\section{Machine Learning Pipeline}
\label{sec:ml}
% ─────────────────────────────────────────────────────────────────────────────

All ML inference is consumed via the HuggingFace Inference API and the z.ai
GLM-5-Turbo API. No local model weights are stored. The Python bridge module
(\texttt{src/python-bridge/client.ts}) wraps all HF calls with error handling
and fallback chains.

\subsection{Models in Production}

\begin{longtable}{@{}p{4.5cm}p{4cm}p{6cm}@{}}
\toprule
\textbf{Model} & \textbf{Task} & \textbf{Input / Output} \\
\midrule
\endhead
\texttt{openai/whisper-large-v3} & Speech-to-text (ASR) & OGG audio bytes → transcribed text + language guess \\
\addlinespace
\texttt{j-hartmann/emotion-english-distilroberta-base} & Text emotion classification & Text string → \{label: anger/joy/sadness/fear/disgust/surprise/neutral, score: 0–1\} \\
\addlinespace
\texttt{audeering/wav2vec2-large-robust-12-ft-emotion-msp-dim} & Dimensional audio emotion & OGG audio bytes → \{valence: 0–1, arousal: 0–1, dominance: 0–1\} \\
\addlinespace
\texttt{papluca/xlm-roberta-base-language-detection} & Language identification & Text string → ISO 639-1 code + confidence \\
\addlinespace
\texttt{HF\_TRANSLATE\_MODEL} (configurable) & Neural machine translation & (source\_text, src\_lang, tgt\_lang) → translated text \\
\addlinespace
\texttt{GLM-5-Turbo} (z.ai, Hermes) & LLM — generation, fallback translation, fallback lang detection & Prompt + context → generated text \\
\bottomrule
\end{longtable}

\subsection{Language Detection — Multi-Stage Confidence Pipeline}

\texttt{detectLanguage()} in \texttt{src/python-bridge/client.ts} implements a
four-stage cascade designed to maximise confidence on Singapore's multilingual
input space (English, Mandarin, Malay, Tamil, Singlish, code-switching):

\begin{lstlisting}[language=]
Stage 1 — Script detection (deterministic, confident: true)
  Unicode range check:
  Tamil:     U+0B80 – U+0BFF
  Malayalam: U+0D00 – U+0D7F
  Devanagari:U+0900 – U+097F  (Hindi)
  Gurmukhi:  U+0A00 – U+0A7F  (Punjabi)
  CJK Unified: U+4E00 – U+9FFF (Chinese/Japanese/Korean)
  → returns {lang, confident: true}

Stage 2 — Malay keyword heuristics (pre-HF, confident: true)
  16-token list: saya, awak, boleh, nak, lah, mah,
                 tidak, ada, ini, itu, dan, dengan,
                 untuk, dari, ke, di
  Threshold: ≥2 keyword hits in the first 8 tokens
  Reason: xlm-roberta consistently mislabels Malay → "ru" (Russian)
  → returns {lang: "ms", confident: true}

Stage 3 — HuggingFace xlm-roberta (ML inference, confident: true)
  Guard: skip if text < 8 chars (model unreliable on short inputs)
  Guard: map output to supported set {en, zh, ms, ta, ml, pa, hi}
  → if in set: returns {lang, confident: true}
  → if out-of-set or API failure: fall to Stage 4

Stage 4 — GLM-5-Turbo LLM fallback (confident: false)
  Prompt: "Detect the language of this text.
           Respond with ONLY an ISO 639-1 code from: en, zh, ms, ta, ml, pa, hi.
           Text: <input>"
  → returns {lang, confident: false}
  Callers receive confident: false and may act on the uncertainty
\end{lstlisting}

When \texttt{confident: false} on the user's \textbf{first message} in a session
(non-English), the bot appends a soft parenthetical to the reply:
\textit{``(Not sure about the language? Just reply in any language and I'll adapt.)''}
This fires at most once per session.

\subsection{Emotion Detection and Distress Scoring}

\subsubsection{Text Emotion Model}

\texttt{j-hartmann/emotion-english-distilroberta-base} classifies text into 7 categories.
The raw output is a confidence score (0–1) for the predicted label — \textbf{not} a
distress magnitude. Post-processing:

\[
\texttt{textBase} = \begin{cases} \texttt{text.score} & \text{if label} \in \{\text{anger, sadness, fear, disgust}\} \\ 0 & \text{otherwise} \end{cases}
\]

The non-distress guard prevents ``thank you!!'' (joy, 0.95) from scoring as a
high-distress message, which would have triggered false escalation.

\subsubsection{Audio Emotion Model (Dimensional)}

The WAV2VEC2 model returns three continuous dimensions:
\begin{itemize}[noitemsep]
  \item \textbf{Valence} (0–1): pleasantness. Low valence = negative affect.
  \item \textbf{Arousal} (0–1): energy/activation. High arousal = excited/loud.
  \item \textbf{Dominance} (0–1): feeling of control (unused in scoring).
\end{itemize}

Audio boost formula:
\[
\texttt{audioBoost} = (1 - \texttt{valence}) \times \texttt{arousal} \times 0.3
\]

\subsubsection{Combined Distress Score}

\[
\texttt{emotion\_score} = \min\!\left((\texttt{textBase} + \texttt{audioBoost}) \times 100,\ 100\right)
\]

\[
\texttt{label} = \begin{cases}
  \text{rage} & \text{if score} > 80 \\
  \text{angry} & \text{if score} > 65 \\
  \text{frustrated} & \text{if score} > 50 \\
  \text{sad} & \text{if score} > 35 \\
  \text{neutral} & \text{otherwise}
\end{cases}
\]

\subsubsection{Emotion Trajectory — Multi-Turn Smoothing}

A single message emotion score is a noisy signal. \texttt{effectiveEmotion()}
in \texttt{src/agents/main/emotionTrajectory.ts} computes a trajectory-smoothed
effective score before tone selection:

\[
\text{recentAvg} = \frac{\sum_{i=1}^{N} w_i \cdot s_i}{\sum_{i=1}^{N} w_i}
\quad \text{where } w_i = 0.6^{N-i} \text{ (recency-decaying)}
\]

\[
\text{effectiveScore} = \max(\text{currentScore},\ \text{recentAvg})
\]

The \texttt{max} ensures the trajectory only \textbf{raises} soothing, never lowers it:
a user who was angry and is now polite still receives warmer responses until the
emotional history fully decays.

\paragraph{Sustained flag.}
\[
\texttt{sustained} = \left|\{s_i : s_i > 65,\ i \in \text{last 3 turns}\}\right| \geq 2
\]

When \texttt{sustained = true}, the tone directive adds an
``ongoing difficulty acknowledgement'' at the start of the reply
(\textit{``I know this hasn't been easy — let's get it sorted''}).

\subsection{Triage Classification}

Before the query agent runs, every message is classified into one of three
categories via regex pattern matching in \texttt{src/agents/triage/classifier.ts}:

\begin{tabular}{@{}lll@{}}
\toprule
\textbf{Category} & \textbf{Type} & \textbf{Action} \\
\midrule
\textbf{Cat 3} — Personal data   & User asks about their own account (balance, payout) & Answer requires officer; escalate \\
\textbf{Cat 2} — Guided topic    & Broad question with a curated guiding set in MongoDB & Interactive Q\&A (guiding questions) \\
\textbf{Cat 1} — Public knowledge & General CPF question answerable from knowledge base & Direct RAG + GLM answer \\
\bottomrule
\end{tabular}

The classifier checks Cat 3 first (personal-data keywords), then Cat 2
(partial-answer patterns), then defaults to Cat 1.

\subsection{Retrieval-Augmented Generation (RAG)}

\begin{lstlisting}[language=]
Query text
  │
  ▼
[1] MongoDB Atlas text index search
    Collection: cpf_knowledge
    Index: { content: "text", title: "text" }
    Returns top-K documents (K = 5 default)
  │
  ▼
[2] Context assembly
    Relevant excerpts concatenated with source metadata
  │
  ▼
[3] GLM-5-Turbo (Hermes, z.ai)
    System prompt: BASE_SYSTEM_PROMPT
      + toneDirective(effectiveEmotion)
      + "Respond in: <lang>"
      + retrieved context (< 4096 tokens)
    User prompt: conversation history + current query
  │
  ▼
[4] Response caching
    Cache key: hash(query + intent + emotion.label + sustained)
    TTL: 5 minutes
    Hit: return cached response without GLM call
\end{lstlisting}

\subsection{Translation Pipeline}

\texttt{translateText()} uses a two-stage fallback:

\begin{enumerate}[noitemsep]
  \item HuggingFace model (\texttt{HF\_TRANSLATE\_MODEL}, configurable): fast, free-tier
  \item GLM-5-Turbo with \texttt{includeSoul: false, disableThinking: true}: utility call,
        avoids the PULSE persona leaking into the translation output
\end{enumerate}

\begin{callout}
\textbf{Critical rule — no double-translation:} The LLM generates replies directly
in the user's detected language (the \texttt{"Respond in: <lang>"} directive). These
native-language replies must \textbf{not} be run through \texttt{translateText()} — that
would double-translate. Translation is used only for fixed UI strings and
officer-to-citizen relay messages.
\end{callout}

\subsection{Background Noise Detection}

For voice messages, a text-to-duration ratio heuristic identifies noisy environments:

\[
\text{noisy} = \left(\frac{|\text{text}|_{\text{chars}}}{\texttt{durationSec}} < 4\right) \land (\texttt{durationSec} \geq 2)
\]

A low characters-per-second ratio indicates Whisper produced sparse output on an
audio clip of meaningful duration — a strong signal for background noise. The user
receives a targeted ``move to a quieter area'' message rather than the generic
``couldn't make out your voice'' message.

% ─────────────────────────────────────────────────────────────────────────────
\section{Escalation and Priority Scoring}
\label{sec:escalation}
% ─────────────────────────────────────────────────────────────────────────────

\subsection{Escalation Analyzer}

\texttt{analyzeEscalation()} in \texttt{src/agents/escalation/analyzer.ts} applies
five ordered layers to decide whether to offer a CPF officer connection:

\begin{lstlisting}[language=]
Layer 1 (highest priority): Explicit request
  Patterns: "speak to a human", "want an officer", "connect me", "escalate"
  → escalate: explicit_request

Layer 2: Personal data
  Patterns: "my CPF balance", "how much do I have", "check my account"
  → escalate: personal_data

Layer 3: Life event
  Patterns: death, disability, divorce, bankruptcy, emigration, estate
  → escalate: life_event

Layer 4: Complaint / dispute
  Patterns: "wrong amount", "not credited", "employer didn't pay"
  → escalate: complaint

Layer 5: Bot failure
  Condition: confidence < 0.4 AND reply contains hedge phrases
  Hedge patterns: "I don't have specific information", "I'm not sure"
  → escalate: bot_failure

Override: emotion_score > 70 → auto-escalate regardless of layers
\end{lstlisting}

\subsection{Priority Scoring Formula}

When a case enters the CCU queue (\texttt{ccu\_queue} MongoDB collection),
a priority score is computed and updated post-escalation:

\[
P = \min\!\left(
  \underbrace{\min(E_{\text{score}}, 95)}_{\text{emotion base}}
  + \underbrace{\min(W \times 0.5,\ 20)}_{\text{wait boost}}
  + \underbrace{15 \cdot \mathbf{1}[\text{financial}]}_{\text{financial flag}}
  + \underbrace{10 \cdot \mathbf{1}[\text{sustained}]}_{\text{sustained flag}},
\ 100\right)
\]

where:
\begin{itemize}[noitemsep]
  \item $E_{\text{score}}$ — emotion distress score (0–100) from \texttt{scoreEmotion}
  \item $W$ — minutes waiting in queue since \texttt{created\_at}
  \item $\text{financial}$ — \texttt{hasFinancialUrgency()} returned true (LIFE\_EVENTS or COMPLAINTS regex fired)
  \item $\text{sustained}$ — \texttt{effectiveEmotion().sustained} is true ($\geq2$ of last 3 turns $>65$)
\end{itemize}

The priority score is refreshed every 30 seconds by the queue refresh timer,
increasing the wait boost as time passes. The dashboard orders cases by descending
priority score.

\subsection{CCU Queue Schema (MongoDB \texttt{ccu\_queue})}

\begin{lstlisting}[language=JavaScript]
{
  queueId:          string,  // UUIDv4
  sessionId:        string,
  userId:           string,  // "tg:<chatId>" | "wa:<phone>"
  display_name:     string,  // Telegram real name or @username
  preferred_lang:   string,  // ISO 639-1
  emotion_score:    number,  // 0–100
  emotion_label:    string,
  summary:          string,  // auto-generated case summary
  query_summary:    string,  // what the user is actually asking
  chat_history:     [
    { role: "user"|"agent"|"officer",
      content: string,
      ts: ISO8601,
      emotion_score?: number,
      emotion_label?: string,
      translated?: string,    // officer message sent to user
      translated_lang?: string }
  ],
  priority_score:   number,  // 0–100, recomputed every 30s
  status:           "waiting" | "assigned" | "resolved",
  assigned_officer: string | null,
  assigned_at:      ISO8601 | null,
  created_at:       ISO8601,
  rating:           number | null  // 1–5 CSAT star rating
}
\end{lstlisting}

% ─────────────────────────────────────────────────────────────────────────────
\section{Messaging Channels}
\label{sec:channels}
% ─────────────────────────────────────────────────────────────────────────────

\subsection{Telegram Bot API}

\begin{tabular}{@{}lp{10cm}@{}}
\toprule
\textbf{Component} & \textbf{Detail} \\
\midrule
Webhook endpoint   & \texttt{POST /webhook/telegram} \\
Registration       & \texttt{setWebhook} with \texttt{allowed\_updates: ["message","edited\_message","callback\_query"]} \\
\texttt{callback\_query} & Required for inline keyboard buttons (escalation offer, CSAT stars). Omitting this field silently breaks the escalation button. \\
Voice messages     & Telegram sends OGG/Opus files. Backend downloads via \texttt{getFile} API. \\
Inline keyboards   & Used for: escalation offer (``Connect to Officer''), CSAT rating (⭐–⭐⭐⭐⭐⭐), officer escalation confirm \\
Rate limits        & \texttt{sendChatAction}: 1 per 5 seconds per chat. Thinking animation uses ~2s interval. \\
\bottomrule
\end{tabular}

\subsection{WhatsApp Meta Cloud API}

\begin{tabular}{@{}lp{10cm}@{}}
\toprule
\textbf{Component} & \textbf{Detail} \\
\midrule
Webhook endpoint   & \texttt{POST /webhook/whatsapp} (JSON, not form-encoded) \\
Hub challenge      & \texttt{GET /webhook/whatsapp} handles Meta's subscription verification (\texttt{hub.challenge} echo) \\
Outbound messages  & Graph API \texttt{POST /{phone\_number\_id}/messages} with \texttt{appsecret\_proof} \\
Signature          & \texttt{X-Hub-Signature-256: sha256=<hmac>} on every inbound webhook \\
Message format     & \texttt{application/json} with \texttt{entry[].changes[].value.messages[]} \\
Voice messages     & WhatsApp sends audio ID; backend fetches via \texttt{GET /{media\_id}} \\
Environment variables & \texttt{WHATSAPP\_PHONE\_NUMBER\_ID}, \texttt{WHATSAPP\_ACCESS\_TOKEN}, \texttt{WHATSAPP\_APP\_SECRET}, \texttt{WHATSAPP\_VERIFY\_TOKEN} \\
\bottomrule
\end{tabular}

\subsection{Channel Abstraction}

Both channels are normalised to a common inbound event shape before hitting \texttt{processInbound()}:
\begin{lstlisting}[language=JavaScript]
{
  userId:   "tg:493251079" | "wa:+6512345678",
  channel:  "tg" | "wa",
  text:     string | null,
  voice:    Buffer | null,  // raw audio bytes
  callbackData: string | null  // inline button payload
}
\end{lstlisting}

% ─────────────────────────────────────────────────────────────────────────────
\section{Officer Dashboard}
\label{sec:dashboard}
% ─────────────────────────────────────────────────────────────────────────────

\subsection{Architecture}

The officer dashboard is a Next.js 14 page served at \texttt{/dashboard}.
It connects to the backend via both REST (initial load) and WebSocket (live updates).

\begin{lstlisting}[language=]
Browser (dashboard page)
  │
  ├── GET /dashboard/queue        ─▶ full queue + stats (initial load)
  ├── GET /dashboard/queue/:id    ─▶ single case with chat history
  ├── POST /dashboard/send/:id    ─▶ officer sends message to citizen
  ├── POST /dashboard/translate   ─▶ on-demand translate citizen message → en
  ├── PATCH /dashboard/resolve/:id ─▶ mark case resolved + trigger CSAT
  └── WSS /dashboard/ws           ─▶ live push events

WebSocket events (server → client):
  { event: "queue_updated",  payload: { queueId, priority_score } }
  { event: "officer_message", payload: { queueId, officerId, ... } }
  { event: "case_resolved",   payload: { queueId } }
  { event: "rating_received", payload: { queueId, rating } }
  { event: "emotion_update",  payload: EmotionEvent }
\end{lstlisting}

\subsection{Real-Time WebSocket Protocol}

WebSocket server is attached to the Express HTTP server instance via \texttt{ws} library
at path \texttt{/dashboard/ws}. The broadcast function fans out to all connected officer
clients:

\begin{lstlisting}[language=JavaScript]
broadcast(event: string, payload: unknown): void {
  const msg = JSON.stringify({ event, payload, ts: Date.now() });
  for (const client of wss.clients) {
    if (client.readyState === WebSocket.OPEN) client.send(msg);
  }
}
\end{lstlisting}

\textbf{Client-side note:} The frontend reads the \texttt{event} field (not \texttt{type}).
The WebSocket connects directly to the backend (\texttt{wss://pulse.nathanielbuilds.cc/dashboard/ws})
bypassing any Next.js proxy that would drop long-lived connections.

\subsection{Officer-to-Citizen Message Flow}

\begin{enumerate}[noitemsep]
  \item Officer types message (English) in dashboard
  \item Backend \texttt{POST /dashboard/send/:queueId}:
    \begin{enumerate}[noitemsep]
      \item Fetches \texttt{preferred\_lang} from queue entry
      \item If \texttt{preferred\_lang $\neq$ en}: calls \texttt{translateText(msg, "en", preferred\_lang)}
      \item Sends translated text via Telegram or WhatsApp adapter
      \item Persists \{role: "officer", content: originalEnglish, translated: translatedText\} to chat history
      \item Broadcasts \texttt{officer\_message} WS event to all dashboard tabs
    \end{enumerate}
\end{enumerate}

This provides a full audit trail: the English original (what the officer intended)
and the translated output (what the citizen received) are both stored.

% ─────────────────────────────────────────────────────────────────────────────
\section{Session Lifecycle and CSAT}
\label{sec:session}
% ─────────────────────────────────────────────────────────────────────────────

Sessions are managed by \texttt{src/services/session/manager.ts}. Each session is
an in-memory object keyed by \texttt{userId}.

\begin{lstlisting}[language=]
Session state:
  conversationHistory: Message[]   (full turn-by-turn history)
  sessionId:           string      (UUIDv4, ties to QueueEntry)
  startedAt:           Date
  lastActivityAt:      Date

Termination conditions:
  1. Citizen signs off ("thank you, bye")  → endSession(userId, "satisfied")
  2. Citizen types /end                    → endSession(userId, "satisfied")
  3. Officer resolves case                 → endSession(userId, "officer", {queueId})
  4. 24h inactivity (sweep timer)          → silent reset, no CSAT

On endSession:
  1. Send CSAT prompt:
     Telegram: inline keyboard ⭐⭐⭐⭐⭐ (callback: "rate:<sessionId>:<1-5>")
     WhatsApp: "Please reply with a number 1–5 to rate this conversation"
  2. Clear conversation history
  3. Reset all pending state (guided Q, pending offer)
\end{lstlisting}

CSAT ratings are stored on the \texttt{QueueEntry} in MongoDB via \texttt{setQueueRating()}.
The dashboard displays the rating once received, and broadcasts a \texttt{rating\_received}
WebSocket event.

% ─────────────────────────────────────────────────────────────────────────────
\section{Ring 4 — Data Layer}
\label{sec:data}
% ─────────────────────────────────────────────────────────────────────────────

\subsection{SQLite — Citizens and Cases}

SQLite (via the \texttt{better-sqlite3} driver) persists:
\begin{itemize}[noitemsep]
  \item \textbf{customers} — user profiles (userId, language preference, vulnerability tier)
  \item \textbf{cases} — service interaction records
  \item \textbf{cpf\_data} — simulated CPF account data (mock, keyed by hash of userId)
\end{itemize}

SQLite is local to the VPS. Resilience is achieved through pm2's automatic restart
on crash (not a multi-node setup; single-server deployment).

\subsection{MongoDB Atlas — Knowledge and Queue}

Active cluster: \texttt{pulse.jvhsxu2.mongodb.net} (Atlas Free Tier, Singapore region).

\begin{tabular}{@{}llp{7cm}@{}}
\toprule
\textbf{Collection} & \textbf{Owned By} & \textbf{Contents} \\
\midrule
\texttt{cpf\_knowledge}         & Hermes / RAG        & CPF scheme documents, FAQ chunks, source metadata \\
\texttt{cpf\_guiding\_questions} & Inbound pipeline   & Curated topic-keyed Q\&A sets for guided flow \\
\texttt{ccu\_queue}             & Escalation pipeline & Live officer queue entries (QueueEntry schema) \\
\texttt{slang}                  & Language agent      & Singapore colloquialisms and glossary terms \\
\bottomrule
\end{tabular}

\subsubsection{Guiding Questions — Topic Key Drift}

The guiding questions collection has a known schema drift between environments:
the committed seed file uses keys \texttt{retirement-income} / \texttt{growing-savings} /
\texttt{home-ownership}, while the live Atlas cluster uses \texttt{retirement} / \texttt{topups} /
\texttt{housing}. Each guiding set carries an \texttt{aliases} array so \texttt{findGuidingSetForQuery()}
matches either vocabulary.

\subsection{In-Memory State}

The following state is in-memory only (lost on backend restart):
\begin{itemize}[noitemsep]
  \item \textbf{UserPrefs} — \texttt{pendingGuiding}, \texttt{pendingOfficerOffer}, \texttt{display\_name}
  \item \textbf{Session history} — conversation turns, CSAT state
  \item \textbf{Queue mirror} — \texttt{queueStore} Map (warmed from MongoDB on startup)
  \item \textbf{Emotion ring buffer} — last 100 emotion events for dashboard
  \item \textbf{Officer status} — connected officers and their availability
  \item \textbf{Session timeout timers} — 24h inactivity sweep handles
\end{itemize}

% ─────────────────────────────────────────────────────────────────────────────
\section{Frontend Architecture}
\label{sec:frontend}
% ─────────────────────────────────────────────────────────────────────────────

\subsection{Technology}

\begin{tabular}{@{}ll@{}}
\toprule
\textbf{Component} & \textbf{Detail} \\
\midrule
Framework       & Next.js 14 (App Router) \\
React version   & React 18 \\
Language        & TypeScript (strict mode) \\
Styling         & Inline \texttt{style=\{\{\}\}} objects + \texttt{globals.css} \\
Fonts           & Fraunces (display/serif) + Bricolage Grotesque (body) via \texttt{next/font/google} \\
Build           & \texttt{next build} → static + dynamic pages; served via \texttt{next start -p 3001} \\
\bottomrule
\end{tabular}

\subsection{Page Structure}

\begin{tabular}{@{}lll@{}}
\toprule
\textbf{Route} & \textbf{Component} & \textbf{Description} \\
\midrule
\texttt{/}               & \texttt{page.tsx}        & Landing page — PULSE pitch, feature cards, nav \\
\texttt{/chat}           & \texttt{chat/page.tsx}   & Citizen chat UI — connects to \texttt{POST /query} \\
\texttt{/cpf}            & \texttt{cpf/page.tsx}    & CPF portal demo (account overview, scheme navigator) \\
\texttt{/dashboard}      & \texttt{CpfDashboard.jsx} & Officer dashboard — live queue, WS, translate, resolve \\
\texttt{/textus}         & \texttt{textus/page.tsx}  & ``Text Us'' demo (conversational entry) \\
\texttt{/correspondence} & stub                     & Correspondence list (scaffold) \\
\texttt{/login}          & stub                     & Singpass login (scaffold) \\
\texttt{/proxy}          & stub                     & Caregiver/proxy portal (scaffold) \\
\texttt{/settings}       & stub                     & User settings (scaffold) \\
\bottomrule
\end{tabular}

\subsection{API Communication}

All frontend API calls go directly to the backend using \texttt{NEXT\_PUBLIC\_BACKEND\_URL}
(empty string in production = same origin via Caddy). The Next.js dev proxy was
bypassed because it drops POST requests after 30 seconds — a known issue with
Next.js 15's fetch proxy implementation. Production requests go same-origin and
never trigger CORS preflight.

% ─────────────────────────────────────────────────────────────────────────────
\section{Deployment and CI/CD}
\label{sec:deploy}
% ─────────────────────────────────────────────────────────────────────────────

\subsection{Process Management (pm2)}

\begin{tabular}{@{}lllll@{}}
\toprule
\textbf{Name} & \textbf{Command} & \textbf{Port} & \textbf{Mode} & \textbf{Restart policy} \\
\midrule
\texttt{pulse-backend}  & \texttt{npx tsx src/gateway/index.ts} & 3000 & fork & On crash, no watch \\
\texttt{pulse-frontend} & \texttt{npx next start -p 3001}       & 3001 & fork & On crash, no watch \\
\bottomrule
\end{tabular}

Both processes run from \texttt{/nat/Project-PULSE}. pm2's persistent process list
means they survive server reboots via \texttt{pm2 startup} + \texttt{pm2 save}.

\subsection{GitHub Webhook Auto-Deploy}

\begin{lstlisting}[language=]
Push to origin/main
  │
  ▼
GitHub sends POST to https://pulse.nathanielbuilds.cc/webhook/github
  │
  ▼
deploy.ts: verify HMAC-SHA256 signature
           check event == "push" and ref == "refs/heads/main"
           respond 202 immediately
           execFile("deploy/deploy-frontend.sh") async
  │
  ▼
deploy-frontend.sh:
  git fetch origin main
  git checkout origin/main -- frontend/   (only frontend/ directory)
  cd frontend && npm install --prefer-offline --no-audit --no-fund
  npm run build                           (next build)
  pm2 restart pulse-frontend
\end{lstlisting}

Only the \texttt{frontend/} directory is checked out from \texttt{main}, preserving any
local backend state. The backend must be restarted manually after backend code changes
(or will pick up changes on next crash-restart since tsx reads \texttt{.ts} files directly).

\subsection{POST — Power-On Self-Test}

The system has a mandatory self-test protocol (\texttt{npm run post}) that must exit 0
before the backend is declared healthy. It verifies:
\begin{itemize}[noitemsep]
  \item Backend HTTP (\texttt{/query})
  \item Frontend HTTP (\texttt{/})
  \item Dashboard WebSocket (\texttt{/dashboard/ws})
  \item Telegram \texttt{getMe} API reachability
  \item Telegram webhook \texttt{allowed\_updates} includes \texttt{callback\_query}
    (auto-repairs if missing — omitting this field silently breaks escalation buttons)
\end{itemize}

% ─────────────────────────────────────────────────────────────────────────────
\section{OpenClaw — AI Agent System}
\label{sec:openclaw}
% ─────────────────────────────────────────────────────────────────────────────

OpenClaw is PULSE's custom multi-agent framework. In the live production system,
the active pipeline runs through the \textbf{main agent} and \textbf{query agent}.
The full domain/language/dialect hierarchy below is the designed target state;
domain agent scaffolds exist in \texttt{src/agents/domain/}.

\subsection{Agent Hierarchy (Designed)}

\begin{lstlisting}[language=]
Orchestrator (routes, never answers directly)
    │
    ├── Domain Agents (deep expertise per domain)
    │   ├── Financial  (CPF, GST, IRAS, tax notices)
    │   ├── Healthcare (Medisave, polyclinic, CHAS, MOH)
    │   ├── Housing    (HDB, rental, conservancy, SERS)
    │   ├── Employment (workpass, MOM, CPF employer)
    │   └── Legal      (fines, court notices, LTA, NEA)
    │
    ├── Language Agents (glossary + contextualisation)
    │   ├── English  (Plain English, Primary 4–6 level)
    │   ├── Chinese  (简体中文 + local colloquialisms)
    │   ├── Malay    (Bahasa Melayu, formal register)
    │   └── Tamil    (தமிழ், formal register + local terms)
    │
    ├── Dialect Agents (12 dialects, first-class)
    │   ├── Hokkien (福建话)    Cantonese (广东话)
    │   ├── Teochew (潮州话)    Hakka (客家话)
    │   ├── Hainanese (海南话)  Bazaar Melayu
    │   ├── Javanese           Singapore Tamil
    │   ├── Spoken Tamil       Malayalam
    │   └── Punjabi            Hindi
    │
    ├── Accessibility Agent
    │   (readability scoring, simplification, TTS, layout)
    │
    └── Guardian Agent
        (scam detection, PII scrubbing, confidence gate)
\end{lstlisting}

\subsection{Agent Communication Protocol}

\begin{lstlisting}[language=JavaScript]
// AgentMessage — inter-agent envelope (src/agents/shared/types.ts)
{
  from:      string,        // sending agent name
  to:        string,        // target agent name
  type:      "request" | "response" | "escalate",
  payload:   AgentPayload,
  context:   RoutingContext,  // userId, sessionId, language, vuln tier, history
  timestamp: ISO8601
}

// RoutingContext
{
  userId:          string,
  sessionId:       string,
  language:        Language,    // "en"|"zh"|"ms"|"ta"
  dialect?:        Dialect,
  vulnerabilityTier: VulnerabilityTier,
  history:         AgentMessage[],
  intent:          DetectedIntent
}
\end{lstlisting}

\subsection{Dynamic Agent Registry}

\texttt{AgentRegistry} (\texttt{src/agents/orchestrator/registry.ts}) maintains
a runtime map of all registered agents. Each agent registers itself at startup
with a capability declaration:

\begin{lstlisting}[language=JavaScript]
{
  name:        string,
  type:        "domain" | "language" | "dialect" | "accessibility" | "guardian",
  domains:     string[],   // which domains this agent handles
  languages:   string[],   // which language codes
  dialects:    string[],   // which dialect codes
  isHealthy(): boolean     // live health check
}
\end{lstlisting}

Adding a new agent requires only: create the agent folder, define the
registration, call \texttt{registry.register(agent)}. Zero existing code changes.

% ─────────────────────────────────────────────────────────────────────────────
\section{Language and Dialect Support}
\label{sec:languages}
% ─────────────────────────────────────────────────────────────────────────────

\subsection{Official Languages}

\begin{tabular}{@{}lll@{}}
\toprule
\textbf{Code} & \textbf{Language} & \textbf{Register / Notes} \\
\midrule
\texttt{en} & English & Plain English, Primary 4–6 reading level \\
\texttt{zh} & Mandarin (Simplified) & 简体中文 + local Singaporean colloquialisms \\
\texttt{ms} & Malay & Bahasa Melayu, formal register \\
\texttt{ta} & Tamil & தமிழ் + local terminology, formal register \\
\bottomrule
\end{tabular}

\subsection{Dialect Coverage (12 dialects)}

\begin{tabular}{@{}lll@{}}
\toprule
\textbf{Code} & \textbf{Dialect} & \textbf{Community Context} \\
\midrule
\texttt{zh-hok} & Hokkien (福建话) & Largest senior dialect group; hawker centres, mature estates \\
\texttt{zh-can} & Cantonese (广东话) & Older generation, HK media consumption \\
\texttt{zh-teo} & Teochew (潮州话) & Concentrated in specific mature HDB estates \\
\texttt{zh-hak} & Hakka (客家话) & Smaller but important for coverage \\
\texttt{zh-hai} & Hainanese (海南话) & Smallest Chinese dialect group \\
\texttt{ms-bms} & Bazaar Melayu & Cross-ethnic colloquial Malay; older generation \\
\texttt{ms-jav} & Javanese & Javanese-descended older Singaporeans \\
\texttt{ta-sin} & Singapore Tamil & Local variety with Malay/English loan words \\
\texttt{ta-spo} & Spoken Tamil & Colloquial; more accessible for non-literate users \\
\texttt{ml}     & Malayalam & Keralite community \\
\texttt{pa}     & Punjabi & Sikh community \\
\texttt{hi}     & Hindi & Broader Indian community, newer arrivals \\
\bottomrule
\end{tabular}

\begin{callout}
\textbf{Why dialects matter:} Many Singaporeans aged 75+ are more comfortable in dialect
than in Mandarin or formal Tamil. The 1979 Speak Mandarin Campaign shifted education away
from dialects, but the current elderly population was already out of school at that time.
For these users, a message in Hokkien or Bazaar Melayu is not a ``translation'' — it is
the \textit{primary language of comprehension}. PULSE treats dialects as first-class
language agents, not localisation afterthoughts.
\end{callout}

% ─────────────────────────────────────────────────────────────────────────────
\section{Service Catalogue (9 Backend Services)}
\label{sec:services}
% ─────────────────────────────────────────────────────────────────────────────

Each service is a self-contained Express Router module following the same internal
structure: \texttt{routes.ts → controller.ts → service.ts → repository.ts}. Dependencies
point inward only; no service imports from another service's internals.

\begin{longtable}{@{}lp{4cm}p{6cm}@{}}
\toprule
\textbf{Service} & \textbf{Owns} & \textbf{Key Responsibilities} \\
\midrule
\endhead
\textbf{Correspondence} & Letters, templates, metadata & Full CRUD with Zod validation; template pattern for all services \\
\textbf{Vulnerability}  & User classification, support tiers & Tier assignment (self-service/guided/high-touch) based on behaviour signals \\
\textbf{Orchestration}  & Channel routing decisions & Determine optimal channel/timing given vulnerability assessment \\
\textbf{Adaptation}     & Language simplification    & Simplify content to target reading level; Flesch-Kincaid scoring \\
\textbf{Delivery}       & Multi-channel dispatch     & SingPost (physical mail), SMS, in-app push, email, voice callbacks \\
\textbf{Notification}   & SMS, push, email reminders & Reminder cadence; 48h digital non-open → SMS follow-up \\
\textbf{Proxy}          & Delegated access           & Caregiver/coordinator relationships; audit trail per action \\
\textbf{Analytics}      & Engagement metrics         & Open rates, drop-off, time-to-completion; feeds vulnerability detector \\
\textbf{Billing}        & Usage metering             & Per-tenant correspondence/voice/mail counts; invoice generation \\
\bottomrule
\end{longtable}

% ─────────────────────────────────────────────────────────────────────────────
\section{Cross-Cutting Concerns}
\label{sec:crosscutting}
% ─────────────────────────────────────────────────────────────────────────────

\subsection{Structured Logging}

All logging uses Pino (\texttt{src/shared/logger.ts}) with structured JSON output:
\begin{lstlisting}[language=JavaScript]
{
  "level": "info",
  "time":  1783762360249,     // Unix ms
  "pid":   2039952,
  "hostname": "localhost",
  "service":  "inbound",      // child logger label
  "userId":   "tg:493251079",
  "msg":      "Inbound message"
  // NEVER: PII, tokens, passwords, full text of user messages
}
\end{lstlisting}

Log levels: \texttt{trace} (verbose dev), \texttt{debug}, \texttt{info},
\texttt{warn} (recoverable), \texttt{error} (needs investigation), \texttt{fatal} (crash).

\subsection{Error Hierarchy}

\texttt{src/shared/errors.ts} defines a typed error hierarchy:
\begin{lstlisting}[language=]
AppError (base)
  ├── NotFoundError      (404)
  ├── UnauthorizedError  (401)
  ├── ForbiddenError     (403)
  ├── ValidationError    (422) — carries Zod issues array
  ├── RateLimitError     (429) — carries retryAfter seconds
  └── ExternalServiceError (502)
        code: "EXTERNAL.<SERVICE>_ERROR"
        isRetryable: boolean
        metadata: any
\end{lstlisting}

The global error handler serialises these to structured JSON. Unhandled errors
return a generic \texttt{500} without leaking stack traces to the client.

\subsection{Message Formatting}

Bot answers follow a government content standard (GOV.UK content design + Singapore SGDS):
\begin{itemize}[noitemsep]
  \item Front-loaded answers: the direct answer comes first
  \item $\leq$25-word sentences
  \item Scannable bullets (capped at 6 per response)
  \item Plain words; one leading emoji (💰 / 🏠 / 🏥 / 📅 / ℹ️)
  \item Always ends with \texttt{cpf.gov.sg} link
\end{itemize}

\texttt{src/shared/formatter.ts} applies deterministic, language-agnostic cleanup:
bullet capping, character-based length cap ($\sim$900 chars, applied before markdown
escaping), CJK-aware sentence boundary detection (\texttt{。！？} in addition to
\texttt{.!?}). Word-count caps are \textbf{never} used — they break on Chinese
(no spaces between characters).

\subsection{Accessibility Standards}

\begin{tabular}{@{}lp{10cm}@{}}
\toprule
\textbf{Standard} & \textbf{Implementation} \\
\midrule
WCAG 2.2 AA       & All UI components. Minimum 4.5:1 contrast ratio enforced. \\
Plain language     & Primary 6 reading level for guided tier; Primary 4 for high-touch \\
Screen readers     & Semantic HTML, \texttt{aria-label} on all interactive elements \\
Responsive layout  & Tested on large-screen phones (common among seniors) and tablets \\
No links in SMS    & Users type known gov.sg URLs manually (scam prevention) \\
\bottomrule
\end{tabular}

% ─────────────────────────────────────────────────────────────────────────────
\section{Vulnerability Detection and Support Tiers}
\label{sec:vulnerability}
% ─────────────────────────────────────────────────────────────────────────────

The system classifies each user into one of three support tiers, dynamically
adjusting the interaction based on detected signals:

\begin{longtable}{@{}lp{5cm}p{5cm}@{}}
\toprule
\textbf{Tier} & \textbf{Trigger Signals} & \textbf{System Behaviour} \\
\midrule
\endhead
\textbf{Self-service}  & Confident digital user; no difficulty signals & Standard UI; full feature access; minimal prompts \\
\textbf{Guided}        & Some difficulty: slow response, repeated corrections, simplified language requested & Simplified language; proactive prompts; step-by-step guidance; larger text hints \\
\textbf{High-touch}    & Repeated failures; age $\geq75$; declared disability; non-digital escalation & Voice assist; human fallback offered proactively; minimal steps; community coordinator alert \\
\bottomrule
\end{longtable}

\textbf{Proactive intervention thresholds:}
\begin{itemize}[noitemsep]
  \item Digital notification unopened after 48h → SMS follow-up with simplified summary
  \item No action taken after 5 days → automated voice-guided callback
  \item Repeated failed interactions detected → flag for community coordinator outreach
\end{itemize}

% ─────────────────────────────────────────────────────────────────────────────
\section{Security Threat Model}
\label{sec:threats}
% ─────────────────────────────────────────────────────────────────────────────

\begin{longtable}{@{}lp{4cm}p{7cm}@{}}
\toprule
\textbf{Threat} & \textbf{Attack Vector} & \textbf{Mitigation} \\
\midrule
\endhead
Webhook spoofing     & Fake Telegram/Meta/GitHub POST & HMAC-SHA256 verification; \texttt{timingSafeEqual} prevents timing oracle \\
Token theft          & JWT in localStorage, XSS steal & httpOnly cookies only; CSP blocks inline scripts \\
Session hijacking    & Cookie theft & Secure; SameSite=Strict; short TTL (15 min access) \\
Replay attack        & Resend old JWT & JTI revocation list in Redis; refresh token rotation \\
SQL injection        & Malicious query input & Parameterised queries via SQLite driver; Zod input validation \\
NoSQL injection      & Malicious MongoDB query & Typed object construction; no \texttt{\$where} queries; Zod sanitisation \\
XSS                  & Script injection in chat & CSP header; no \texttt{dangerouslySetInnerHTML}; markdown escaping in formatter \\
Clickjacking         & iframe embedding & \texttt{X-Frame-Options: DENY} \\
DDoS                 & Traffic flood & Rate limiting per IP/user; Caddy connection limits \\
Scam impersonation   & Fake gov messages to users & No links in SMS/email; verifiable reference codes; scam education in flow \\
PII leakage via logs & Log PII by accident & Pino child loggers never log message content; userId only (hashed/opaque) \\
MITM                 & HTTP traffic interception & TLS enforced; HSTS \texttt{max-age=31536000}; HTTP 301 redirect to HTTPS \\
Credential exposure  & Hardcoded secrets & All secrets in \texttt{.env} (gitignored); \texttt{.env.example} commits only \\
\bottomrule
\end{longtable}

% ─────────────────────────────────────────────────────────────────────────────
\section{Aspirational Infrastructure (Not Yet Provisioned)}
\label{sec:aspirational}
% ─────────────────────────────────────────────────────────────────────────────

The following components are documented in the architecture and scaffolded in code
but have not been provisioned in the current single-VPS deployment:

\begin{tabular}{@{}lll@{}}
\toprule
\textbf{Component} & \textbf{Planned Role} & \textbf{Status} \\
\midrule
Apache Kafka       & Async inter-service event bus & Deps removed (kafkajs was unused) \\
Redis              & Rate limit storage, JWT revocation, response cache & In-memory shim used instead \\
PostgreSQL         & Schema-per-service persistence & SQLite used instead \\
AWS/Azure Gov Cloud & Multi-AZ production hosting & Single VPS (pulse.nathanielbuilds.cc) \\
Docker + Kubernetes & Container orchestration & pm2 used instead \\
CloudFront / Cloudflare CDN & Edge caching, WAF, DDoS shield & Caddy on VPS \\
Singpass IdP integration & Citizen authentication & JWT dev-mode stub \\
SingPost API        & Physical mail dispatch & Adapter scaffold in place \\
OpenTelemetry/Grafana & Distributed tracing, metrics & Pino structured logs only \\
LaunchDarkly/Unleash & Feature flags & Direct code deployment \\
\bottomrule
\end{tabular}

% ─────────────────────────────────────────────────────────────────────────────
\section{Environment Variables Reference}
\label{sec:env}
% ─────────────────────────────────────────────────────────────────────────────

\begin{longtable}{@{}lp{5cm}l@{}}
\toprule
\textbf{Variable} & \textbf{Purpose} & \textbf{Required} \\
\midrule
\endhead
\texttt{TELEGRAM\_BOT\_TOKEN}       & Telegram Bot API authentication & Yes \\
\texttt{TELEGRAM\_WEBHOOK\_SECRET}  & Telegram webhook token verification & Recommended \\
\texttt{WHATSAPP\_PHONE\_NUMBER\_ID} & Meta Cloud API phone number ID & Yes (WA) \\
\texttt{WHATSAPP\_ACCESS\_TOKEN}    & Meta Cloud API bearer token & Yes (WA) \\
\texttt{WHATSAPP\_APP\_SECRET}      & HMAC-SHA256 webhook signing key & Yes (WA) \\
\texttt{WHATSAPP\_VERIFY\_TOKEN}    & Meta hub-challenge verification & Yes (WA) \\
\texttt{GITHUB\_WEBHOOK\_SECRET}    & GitHub auto-deploy HMAC key & Yes (auto-deploy) \\
\texttt{MONGODB\_URI}               & Atlas cluster connection string & Yes \\
\texttt{HF\_API\_KEY}               & HuggingFace Inference API key & Yes \\
\texttt{HF\_TRANSLATE\_MODEL}       & Translation model endpoint & Optional (GLM fallback) \\
\texttt{GLM\_API\_KEY}              & z.ai GLM-5-Turbo API key & Yes \\
\texttt{GLM\_BASE\_URL}             & z.ai API base URL & Yes \\
\texttt{JWT\_SECRET}                & JWT signing secret & Yes \\
\texttt{NEXT\_PUBLIC\_BACKEND\_URL} & Frontend API base (empty = same-origin) & Frontend \\
\texttt{NODE\_ENV}                  & \texttt{production} / \texttt{development} & Yes \\
\bottomrule
\end{longtable}

% ─────────────────────────────────────────────────────────────────────────────
\section{Key Data Flows — End-to-End Sequence}
\label{sec:flows}
% ─────────────────────────────────────────────────────────────────────────────

\subsection{Happy Path: Citizen Asks a CPF Question via Telegram}

\begin{lstlisting}[language=]
Citizen (Telegram)       Backend                     External APIs
      │                     │                              │
      │── "How much CPF ──▶ │                              │
      │    can I withdraw?" │                              │
      │                     │── sendChatAction(typing) ──▶ Telegram
      │                     │                              │
      │                     ├── detectEmotion() ──────────▶ HF (distilroberta)
      │                     ├── detectLanguage() ─────────▶ HF (xlm-roberta) or GLM
      │                     │                              │
      │                     │── classifyQuery() ──▶ Cat 1 (public_knowledge)
      │                     │                              │
      │                     │── searchKnowledge() ─────── MongoDB Atlas (cpf_knowledge)
      │                     │                              │
      │                     │── chatComplete() ───────────▶ GLM-5-Turbo (z.ai)
      │                     │                              │
      │                     │── formatReply() ──▶ (local, no API)
      │                     │                              │
      │                     │── analyzeEscalation() ──▶ (local regex)
      │                     │                              │
      │◀── "CPF members can ──│── sendMessage() ───────────▶ Telegram Bot API
      │    withdraw..." msg   │                              │
      │                     │── broadcast(emotion_update) ▶ WS (officer dashboard)
\end{lstlisting}

\subsection{Escalation Path: Citizen Escalated to Officer}

\begin{lstlisting}[language=]
Citizen (Telegram)     Backend               Officer (Dashboard)   MongoDB Atlas
      │                   │                          │                  │
      │─ "I want to speak  ─▶│                          │                  │
      │   to a real person" │                          │                  │
      │                   │── analyzeEscalation() → explicit_request     │
      │                   │── postToQueue() ──────────────────────────── ccu_queue
      │◀─ inline keyboard ──│                          │                  │
      │   [Connect to Officer]                         │                  │
      │── (tap button) ──▶│── callback_query            │                  │
      │                   │── doEscalate()             │                  │
      │                   │── hasFinancialUrgency() ──▶ (local regex)     │
      │                   │── computePriorityScore() ─▶ (local formula)   │
      │                   │── updateQueuePriority() ───────────────────── ccu_queue
      │                   │── broadcast(queue_updated) ──▶ WS           │
      │◀─ "Connecting you  │                          │                  │
      │   to an officer..." │                          │                  │
      │                   │          [officer opens dashboard]            │
      │                   │◀─ GET /dashboard/queue ──── officer fetches  │
      │                   │                          │                  │
      │                   │         [officer types reply]                 │
      │                   │◀─ POST /dashboard/send ──────────────────────│
      │                   │── translateText("en"→"ms")──▶ GLM/HF        │
      │◀─ (translated msg)─│── sendMessage() ──────────▶ Telegram        │
      │                   │── broadcast(officer_message)──▶ WS           │
      │                   │── appendToQueueHistory()─────────────────── ccu_queue
\end{lstlisting}

% ─────────────────────────────────────────────────────────────────────────────
\section{Testing and Self-Verification}
\label{sec:testing}
% ─────────────────────────────────────────────────────────────────────────────

\begin{tabular}{@{}lll@{}}
\toprule
\textbf{Test Type} & \textbf{Command} & \textbf{Scope} \\
\midrule
Unit tests           & \texttt{npm test}      & 107 tests — emotion scoring, trajectory, formatter, flow predicates \\
POST self-test       & \texttt{npm run post}  & Live system health: HTTP, WS, Telegram API, webhook config \\
TypeScript typecheck & \texttt{npx tsc --noEmit} & Zero errors enforced before push \\
Frontend build       & \texttt{npm run build} (frontend/) & All 12 pages prerender without errors \\
\bottomrule
\end{tabular}

Key test files:
\begin{itemize}[noitemsep]
  \item \texttt{src/agents/main/emotionTrajectory.test.ts} — rising/falling/spike/cold-start/sustained cases
  \item \texttt{src/agents/main/emotion.test.ts} — distress label boundary conditions
  \item \texttt{src/agents/triage/classifier.test.ts} — Cat 1/2/3 classification
  \item \texttt{src/gateway/inbound.flow.test.ts} — 12-case characterisation harness
  \item \texttt{src/shared/formatter.test.ts} — CJK boundary, bullet cap, URL preservation
\end{itemize}

% ─────────────────────────────────────────────────────────────────────────────
\section{Project Roadmap}
\label{sec:roadmap}
% ─────────────────────────────────────────────────────────────────────────────

\begin{longtable}{@{}lp{8cm}l@{}}
\toprule
\textbf{Item} & \textbf{Description} & \textbf{Priority} \\
\midrule
\endhead
History persistence (B9) & Conversation history lost on backend restart; move to MongoDB & High \\
Singpass integration & Full OAuth2 flow with NRIC verification & High \\
Domain agents (5) & Implement Financial, Housing, Healthcare, Employment, Legal agents with curated knowledge bases & High \\
Dialect agents (12) & Full dialect response generation + glossary lookup & Medium \\
PostgreSQL provisioning & Replace SQLite with schema-per-service PostgreSQL & Medium \\
Redis caching & Replace in-memory rate limiter and response cache with Redis & Medium \\
Physical mail (SingPost) & Complete delivery adapter for postal dispatch & Medium \\
Voice IVR & Twilio/AWS Connect integration for phone channel & Medium \\
Kafka event bus & Async inter-service communication via topics & Low \\
Kubernetes + Docker & Container orchestration for horizontal scaling & Low \\
\bottomrule
\end{longtable}

% ─────────────────────────────────────────────────────────────────────────────
\appendix
\section{Glossary}
% ─────────────────────────────────────────────────────────────────────────────

\begin{tabular}{@{}lp{11cm}@{}}
\toprule
\textbf{Term} & \textbf{Definition} \\
\midrule
OpenClaw  & PULSE's custom multi-agent AI framework (name is internal) \\
Hermes    & The LLM component — GLM-5-Turbo on z.ai, used for RAG generation + translation fallback \\
CCU       & Contact Centre Unit — the officer team that handles escalated cases \\
RAG       & Retrieval-Augmented Generation — combining document retrieval with LLM generation \\
HF        & HuggingFace — ML model inference provider \\
CSAT      & Customer Satisfaction Score — 1–5 star rating collected at session end \\
triage    & Three-category query classifier (Cat 1 public / Cat 2 guided / Cat 3 personal) \\
tsx       & TypeScript execute — runs \texttt{.ts} files directly without a compile step \\
pm2       & Node.js process manager — keeps processes alive, logs, restarts on crash \\
Caddy     & Reverse proxy + automatic TLS; routes \texttt{/api/*} → :3000, frontend → :3001 \\
GLM       & General Language Model — z.ai's GLM-5-Turbo API (Hermes uses this) \\
JTI       & JWT ID — unique token identifier used for revocation \\
HMAC      & Hash-based Message Authentication Code — used for webhook signature verification \\
OGG/Opus  & Audio codec used by Telegram voice messages \\
\bottomrule
\end{tabular}

% ─────────────────────────────────────────────────────────────────────────────
\vfill
\begin{center}
\small\textcolor{gray}{
  PULSE Framework — System Architecture Document \\
  Version 3.0 — July 2026 \\
  \texttt{OriginalNneo/Project-PULSE} — branch \texttt{main} \\[4pt]
  This document reflects the live production system as of the date above. \\
  The aspirational infrastructure described in Section~\ref{sec:aspirational} \\
  represents planned but not yet provisioned components.
}
\end{center}
% ─────────────────────────────────────────────────────────────────────────────

\end{document}
